\chapter{Code}\label{app:code}
The implementation is available at \url{https://github.com/pmccarthy3901/posesolving}. 
Table \ref{tab:code_map} maps the sections of the thesis to the corresponding Python modules.


\begin{table}
  \centering 
  \begin{tabular}{lll}
    \toprule 
    Section & Module & Contents \\
    \midrule 
    \ref{sec:anser_sim} & \lstinline[style=python] |anser/| directory & Anser forward model simulation. \\
    \ref{sec:dataset} & \lstinline[style=python] |data/build_dataset.py| & Dataset generation. \\
    \ref{sec:error} & \lstinline[style=python] |models/train.py| & Error metrics. \\
    \ref{sec:methods-lm} & \lstinline[style=python] |models/LM_solve.py| & LM implementation. \\
    \ref{sec:network} & \lstinline[style=python]|models/FFNN_network.py| & NN configuration. \\
    \ref{sec:methods-loss} & \lstinline[style=python] |models/train.py| & Loss function. \\
    \ref{sec:nn_training} & \lstinline[style=python] |models/train.py| & Network training. \\
    \ref{sec:hybrid} & \lstinline[style=python] |data/reparametrisation.py| & Reparametrisation. \\
    \ref{sec:results-lm} & \lstinline[style=python] |demos/basin_sweep.ipynb| & Basin sweep. \\
    \ref{sec:three_way} & \lstinline[style=python] |demos/solver_demos.ipynb| & Three-way comparison. \\
    \bottomrule 
  \end{tabular}
  \caption{Mapping between thesis sections and Python modules.}
  \label{tab:code_map}
\end{table}

